\section{読み方}
\label{sec:reading}

本書の読者は、DST の多体計算を書くプログラマである。物理の総説をここで再構成しない。実装が合わせる相手は、リポジトリの Lean モジュールである。モジュール文書と定理が仕様であり、論文の段落は索引である。

\subsection{三つの地位}

関数を置く前に、その関数の地位をファイル先頭のコメントへ書く。

\begin{description}
\item[\dstrec] Lean にある定義または定理と、引数の意味まで一致させる。名前は Lean の識別子をそのまま使う。
\item[\dstwork] 動態の論文が「書かれた作用ではない」と分けている近傍系である。識別子に \texttt{working} を付ける。記録用のソルバから呼ばない。
\item[\dstopen] Lean に積分器も接触則も無い界面である。プログラムが値を埋めるなら、呼び出し側が渡した仮定としてログへ残す。論文の運動方程式へ昇格させない。
\end{description}

地位の違う式を一つの更新関数に混ぜると、テストがどの主張を見ているか分からなくなる。モジュール境界は第\ref{sec:modules}章で固定する。

\subsection{二つのトラック}

DST には、少なくとも次の二つの計算がある。

\begin{enumerate}
\item 粒子の内在角 \((\phi_a,\theta_a)\) の時間発展。書かれた作用の閉形式は第\ref{sec:written}章である。
\item 共有テトラッドのチャート上の外部軌道。水星の近日点余剰は、外部の軌道要素を動径回帰へ入れた量である。太陽光線の偏折は、外部のヌル軌道の方位角余剰である。
\end{enumerate}

チャートの数値と、内在角の閉形式の出力は、別のテストで見る。対応は第\ref{sec:checks}章に書く。層の節に出る \(\alpha,\beta\) は干渉因子 \(\gamma_s\) の引数であり、粒子状態のラピディティとは別の読みである。層力は粒子の内在角を読まない。

\subsection{多体という言葉}

Lean が多体について固定しているのは、有限集合上の静的な力和である。時間ステップは取らない。粒子数を変える反応も置かない。時刻を進める一段が必要なら、第\ref{sec:contract}章の仮定をログに残したうえで、作業模型としてだけ実行する。
